\section{Discussion}
\label{sec:discussion}

\paragraph{Semantic-addressability principle.}
The strongest evidence is not that a LUT can reproduce arithmetic, but that
correct Q/K state--response alignment predicts held-out responses beyond
global, coarse, random, and shuffled controls. Binary state alone is therefore
insufficient; the address must preserve architecture-specific semantics. The
mixed-radix tuple has no hash collision, yet it can still alias distinct
continuous contexts. Conditional variance and shuffled controls quantify this
remaining ambiguity rather than assuming a binary code is automatically a
meaningful key.

\paragraph{Dynamic-contract principle.}
Local reconstruction and downstream fidelity are different requirements.
Static prototypes can preserve a response mode while shifting its
time/channel moments. BN and LIF thresholds can amplify that shift, especially
over multiple time steps. TCSLU addresses this second requirement by restoring
the statistics of the current presented to the stateful path. In this sense,
semantic addressing determines \emph{what} to retrieve and dynamic calibration
determines whether the retrieved response remains compatible with spiking
dynamics.

\paragraph{Why explicit backoff matters.}
Unseen and low-support addresses are unavoidable in a finite calibration set.
Returning an unsupported prototype would turn uncertainty into an unreported
error. Global-residual lookup instead defines a valid coarse prediction, and
the hierarchy records how often finer semantics are trusted. E5 shows that
this mechanism is reliable but not sufficiently compact by itself; E7 is
needed to retain semantic gain under the entry budget.

\paragraph{Reconstruction is not classification.}
The E1 and E7 results establish local conditional structure, not accuracy
improvement. End-to-end replacement is therefore evaluated as a separate gate.
The calibrated rows show that local structure can be used without materially
damaging the frozen classifier in the tested QKFormer settings. Small positive
accuracy differences are treated as no-loss fluctuations. This separation is
important because historical trainable adapters could adapt around broken
address assignments and therefore did not provide a clean semantic test.

\paragraph{Architecture-conditioned boundary.}
Unmodified Spikformer remains far from lossless replacement despite
temporal/channel calibration. After changing its SSA interaction to a
QKFormer-like $\operatorname{gate}(Q)\odot K$ contract, the replacement gap
shrinks substantially. This intervention supports lookup compatibility as an
architectural property. It does not show that fine Q/K addresses always
dominate coarse state: in the modified Spikformer, aligned lookup beats
shuffling but is effectively tied with token/channel. The present principle is
therefore bounded to QK-compatible interactions, not all spiking neurons.

\paragraph{Event-data implementation boundary.}
The CIFAR10-DVS study contains two valid LUT-only deployments. Route 1 reaches
$83.9\%$ and remains the strongest result in the fixed route family. Route 3
reaches $81.7\%$ from an $81.2\%$ from-scratch SpiLiFormer teacher while
bypassing the original projection, showing that the replacement mechanism can
execute on a second architecture. The official $86.7\%$ clean result was not
reproduced, however, and neither route reaches the registered $84.0\%$ gate.
Route 3 is therefore bounded cross-architecture implementation evidence, not a
strong absolute-accuracy or broad-generalization result.

\paragraph{Analytical cost versus measured systems behavior.}
Supported-entry and metadata-aware byte counts make the table accounting more
honest, but they still omit cache behavior, alignment, address generation,
control logic, and memory traffic. Likewise, grouped and native operator
audits establish numerical and classification fidelity for specified operator
classes, not a hardware speedup. Claims about latency, SRAM, area, or energy
require an implementation that removes the original computation and measures
the resulting system.
